\printedpage{24}
% head: A. Grothendieck : Dérivateurs, Ch. XVII, Édition des Universités de Montpellier II et Paris VII

\msp{45}in $\mathcal{M}(I^\circ)$, let us consider the triples
\[
(J,\ Z,\ i,\ p)\ ,
\]
where $J\subset I$ is an open set, $Z$ an object of $\mathcal{M}(J^\circ)$, and $i$, $p$ are arrows
\[
\left\{
\begin{aligned}
& X|J \xrightarrow{\ i\ } Z \xleftarrow{\ p\ } Y|J \\
& \text{with } i\in C_{J^\circ},\ p\in W_{J^\circ}.
\end{aligned}
\right.
\]
We order the quadruples by extension. If we have a totally ordered set, of $(J_\alpha,Z_\alpha,i_\alpha,p_\alpha)$, putting $J=\bigcup I_\alpha$, we find unique $Z$, $i$, $p$ on $J$ which extend % [sic] \bigcup I_\alpha -> \bigcup J_\alpha, the union is of the open sets J_\alpha
the given $Z_\alpha$, $i_\alpha$, $p_\alpha$, moreover we will have $i\in C_{J^\circ}$, $p\in W_{J^\circ}$, since membership in $C_{J^\circ}$ resp. $W_{J^\circ}$ is a local question on $J$. So by Zorn, we are reduced to proving that if $(J,Z,i,p)$ is maximal, we have $J=I$ – whence the desired factorization $f=pi$. So to prove that if % [sic] d'ou -> d'où, missing accent
$J\neq I$, the quadruple is not maximal. Let us consider the nonempty closed set $I\backslash J$. $I$ is assumed artinian, there is a \textit{minimal element} $i_0$ of $I\backslash J$. Let us put
\msp{46}
\[
J' \;=\; J\cup\{i_0\} \;=\; J\cup I/i_0,
\]
and let us show that we can extend the triple $Z$, $i$, $p$ to $Z'$, $j'$, $p'$ on $J'$. Putting $K'=I/i_0$, % [sic] j' -> i', the triplet is Z, i, p
$K=J\cap K=I_{<i_0}$, we thus have % [sic] J\cap K -> J\cap K', as in the diagram below
\[
\begin{tikzcd}
K=J\cap K' \arrow[r, hook] \arrow[d, hook] & K'=I/i_0 \arrow[d, hook] \\
J \arrow[r, hook] & J'=J\cup K'.
\end{tikzcd}
\]

We then see that the extensions of $(Z,i,p)$ from $J$ to $J'$ correspond to those of $(Z|K',i|K',p|K')$ to $K$. For this, we note that for an arrow $\varphi$ between objects of $\mathcal{M}(J'^\circ)$, this arrow is in $C_{J'^\circ}$ resp. $W_{J'^\circ}$, if and only if its restrictions to $J$ and to $K'$ are.

We are thus reduced to the standard case of an ordered set with final object $e_0$, an ordered set which we will denote by $I$, and the ordered set $I'$ of objects $i<e_0$, i.e. $I'=I\backslash\{i_0\}$. % [sic] I\backslash\{i_0\} -> I\backslash\{e_0\}, I' consists of the i < e_0
We have $f:X_\bullet\longrightarrow Y_\bullet$, an arrow in $\mathcal{M}(I^\circ)$, a factorization of $f|I'$ as $p'\circ i'$ with $i'\in C_{I'^\circ}$, $p'\in W_{I^\circ}$, % [sic] W_{I^\circ} -> W_{I'^\circ}, factorisation of f|I' on I'
\msp{47}
and it is a matter of extending it to a similar factorization of $f$. We note that giving an object $X_\bullet$ of $\mathcal{M}(I^\circ)$ is equivalent to giving an object $X'_\bullet$ of $\mathcal{M}(I'^\circ)$ (the restriction of $X$ to $I'$) and an arrow
\[
X'_{I'} \xrightarrow{\ \alpha_X\ } X_{e_0}=X_0
\]
in $\mathcal{M}$, where $X'_{I'}=\varinjlim\limits_{i'\in I'}X'_{i'}$. Thus giving the arrow $f:X_\bullet\longrightarrow Y_\bullet$ amounts to giving $f'=f|I':X'_\bullet\longrightarrow Y'_\bullet$, and an arrow $f_0:X_0\longrightarrow Y_0$ making the square commutative
\printedpage{25}
% head: A. Grothendieck : Dérivateurs, Ch. XVII, Édition des Universités de Montpellier II et Paris VII
\[
\begin{tikzcd}
X'_{I'} \arrow[r, "\alpha_X"] \arrow[d, "f'_{I'}"'] & X_0 \arrow[d, "f_0"] \\
Y'_{I'} \arrow[r, "\alpha_Y"] & Y_0\ .
\end{tikzcd}
\]

Likewise, finding a factorization of $f$ as $pi$, which extends a given factorization $f'=p'i'$ on $J'$, amounts to giving a commutative diagram
\[
\begin{tikzcd}[column sep=small, row sep=small]
& X'_{I'} \arrow[rr, "\alpha_X"] \arrow[dl, "i'_{I'}"'] \arrow[dd, "f'_{I'}"] & & X_0 \arrow[dl, "i_0"'] \arrow[dd, "f_0"] \\
Z_{I'} \arrow[rr, "\alpha_Z"'] \arrow[dr, "p'_{I'}"'] & & Z_0 \arrow[dr, "p_0"'] & \\
& Y'_{I'} \arrow[rr, "\alpha_Y"'] & & Y_0\ ,
\end{tikzcd}
\tag*{$(*)$}
\]
\msp{48 vide}
\msp{49}
where the object $Z_0$, and the three arrows $\alpha_Z$, $i_0$, $p_0$ are to be determined. If we have $i'\in C_{I'^\circ}$, $p\in W_{I'^\circ}$, % [sic] p -> p', the given factorisation is f'=p'i'
then we conclude by theorem 2.1. $1^\circ$ (cor. 2.2) that
\[
i'_{I'}\ \in\ C\ ,
\]
and if we want $i\in C_{I^\circ}$, $p\in W_{I^\circ}$, it only remains to require
\[
p_0\in W, \qquad (i'_{I',I}:\underbrace{Z_{I'}\amalg_{X_{I'}}X_0}_{Z'_0}\longrightarrow Z_0)\in C,
\]
by theorem 2.1. $2^\circ$. We will then construct the diagram $(*)$, from the similar diagram below
\[
\begin{tikzcd}
& X_{I'} \arrow[rr, "\alpha_X"] \arrow[dl, "i'_{I'}\in C"'] \arrow[dd, "\text{cocart.}" {pos=0.25}, "f'_{I'}" {pos=0.75}] & & X_0 \arrow[dl, "i'_0\in C"'] \arrow[dd, "f_0" {pos=0.25}] \arrow[dr, dotted, "i_0=j_0i'_0"] & \\
Z_{I'} \arrow[rr, "\alpha'_Z"] \arrow[dr, "p'_{I'}"'] & & Z'_0=Z_{I',I} \arrow[rr, dotted, "j_0\in C"] \arrow[dr, "p'_0"'] & & Z_0 \arrow[dl, dotted, "p_0\in W"] \\
& Y_{I'} \arrow[rr, "\alpha_Y"] & & Y_0 &
\end{tikzcd}
\]
(with $Z_0$ of $(*)$ replaced by $Z'_0=Z_{I',I}$) by factoring the canonical arrow $p'_0:Z'_0\longrightarrow Y_0$ as
\[
p'_0 \;=\; p_0j_0\ , \qquad \text{with } p_0\in W,\ j_0\in C\ ,
\]
and putting
\[
i_0 \;=\; j_0i'_0\ ,
\]
which completes the construction.

We have thus proved the

\printedpage{26}
% head: A. Grothendieck : Dérivateurs, Ch. XVII, Édition des Universités de Montpellier II et Paris VII

\msp{50}\textbf{Theorem 3.1.} \textit{Let $(\mathcal{M},W,C)$ be a category with cofibrations (cf. page \textup{39}), $\mathcal{M}$ stable under finite $\varinjlim$. For any ordered set $I$, let us consider the corresponding tuple $(\mathcal{M}(I^\circ),W_{I^\circ},C_{I^\circ})$, where $C_{I^\circ}$ is defined in paragraph \textup{2} (cf. on this point theorem \textup{2.1}). Suppose that $I$ is an \textup{artinian} ordered set (e.g. $I$ locally finite, i.e. the $I/i$ finite). Suppose moreover that we are in one of the following two cases:} % [sic] odonné -> ordonné, missing letter
\begin{itemize}[leftmargin=2.5em]
\item[a)] \textit{The space \textup{Esp}$(I)$ is noetherian, i.e. every increasing sequence of open sets of $I$ is stationary (e.g. $I$ finite).}
\item[b)] \textit{$C$ and $TC=C\cap W$ are stable, not only under ordinary composition, but also under transfinite composition. In this case, we also assume $\mathcal{M}$ stable under small filtering $\varinjlim$ (hence under arbitrary small $\varinjlim$).}
\end{itemize}
\textit{Under these conditions, $(\mathcal{M}(I^\circ),W_{I^\circ},C_{I^\circ})$ is also a category with cofibrations.}

\textbf{Proposition 3.2.} \textit{Let $u:I\longrightarrow I'$ be an increasing map between ordered sets, whence a pair of adjoint functors}
\[
\begin{tikzcd}
\mathcal{M}(I^\circ) \arrow[r, shift left, "u_{!}"] & \mathcal{M}(I'^\circ) \arrow[l, shift left, "u^*"]
\end{tikzcd}\ .
\tag*{(3.2.1)}
\]
\msp{51}
\begin{itemize}[leftmargin=2.5em]
\item[a)] \textit{The functor $u^*$ is compatible with the localizers $W_{I'^\circ}$, $W_{I^\circ}$. If $I$ is discrete, or $u$ étale, then $u^*$ maps $C_{I'^\circ}$ into $C_{I^\circ}$. The same is still true when $u$ is locally cofiltering, i.e. if for every $i'\in\operatorname{Ob}I'$, the category $i'\backslash I$ is locally filtering, i.e. its connected components are cofiltering (condition $\mathrm{PS}_{\mathrm{g}}$ \textup{1} on $i'\backslash I$)} \textsuperscript{(}\footnote[16]{no need to speak of ``weakly'' locally cofiltering, since these are ordered sets where there is no need to distinguish between the strong and weak notions.}\textsuperscript{)}.
\item[b)] \textit{In all cases, $u_!$ transforms $C_{I^\circ}$ into $C_{I'^\circ}$, $TC_{I^\circ}$ into $TC_{I'^\circ}$, finally $\mathcal{M}(I^\circ)_{\mathrm{c}}$ into $\mathcal{M}(I'^\circ)_{\mathrm{c}}$.}
\end{itemize}

\textsc{Proof}. Let us prove a), in the most awkward case, where we assume only that $u$ is locally weakly cofiltering. We apply II 7.4 (p. 92), which tells us that $\forall i\in\operatorname{Ob}I$,
\msp{52}
considering the morphism induced by $u$
\[
I/i \longrightarrow I'/i' \qquad (\text{where } i'=u(i)),
\]
this is \textit{(globally)} weakly cofiltering, and that for every open set $U$ of $I/i$, denoting by $U'$ the open set generated by its image, the induced morphism $v:U\longrightarrow U'$ is still weakly cofiltering, and a fortiori \textit{cofinal.} Let then $f':X'\longrightarrow Y'$ be an arrow in $C_{I'^\circ}$, let us prove that
\[
\underbrace{u^*(f')}_{f}\ :\ \underbrace{u^*(X')}_{X}\ \longrightarrow\ \underbrace{u^*(Y')}_{Y}
\]
is in $C_{I^\circ}$, i.e. that for every $i\in I$, and $U$ an open set of $I/i$, in the canonical diagram
\printedpage{27}
% head: A. Grothendieck : Dérivateurs, Ch. XVII, Édition des Universités de Montpellier II et Paris VII
\[
\begin{tikzcd}
X_U \arrow[r] \arrow[d] \arrow[dr, phantom, "\text{cocart.}"] & X(i) \arrow[d] \\
Y_U \arrow[r] & Y_{U,i} \arrow[d, "v_{U,i}"] \\
& Y(i)\ ,
\end{tikzcd}
\tag*{$(*)$}
\]
\msp{53}
the arrow $v_{U,i}$ is in $C$. The diagram $(*)$ maps into the similar diagram for $f':X'\longrightarrow Y'$, $i'=u(i)$, and the open set $U'\subset I'/i'$ generated by $u(U)$:
\[
\begin{tikzcd}[column sep=small, row sep=small]
X_U \arrow[rr] \arrow[dr, "f_U"] \arrow[dd, "c_U^{X'}"'] & & X(i) \arrow[dr] \arrow[dd, "\text{cocart.}"' {pos=0.2}, "\mathrm{id}"' {pos=0.75}] & & \\
& Y_U \arrow[rr] \arrow[dd, "c_U^{Y'}"' {pos=0.3}] & & Y_{U,i} \arrow[dr, "v_{U,i}"] \arrow[dd, "c^{f'}_{U,i}"] & \\
X'_{U'} \arrow[rr, "\text{cocart.}"' {pos=0.6}] \arrow[dr, "f'_{U'}"'] & & X'(i') \arrow[dr] & & Y(i) \arrow[dd, "\mathrm{id}"] \\
& Y'_{U'} \arrow[rr] & & Y'_{U',i'} \arrow[dr, "v_{U',i'}"'] & \\
& & & & Y'(i')\ .
\end{tikzcd}
\tag*{$(**)$}
\]

By cofinality, the two vertical arrows $c_U^{X'}$ and $c_U^{Y'}$ are isomorphisms, in addition to $X(i)\xrightarrow{\ id\ }X'(i')$, so the fourth arrow of the cube, $c^{f'}_{U,i}$, is an isomorphism. Since $v_{U',i'}$ is in $C$, $v_{U,i}$ too, we win.

Let us prove b). So let
\[
f\ :\ X\longrightarrow Y
\]
be an arrow in $C_I$, and let
\[
f' \;=\; u_!(f)\ :\ u_!(X)=X' \longrightarrow u_!(Y)=Y',
\]
let us prove that it is in $C_{I'^\circ}$, hence that
\msp{54}
for two open sets $U'\subset V'$ of $I'$ (where we can assume that $V'=I'/i'$), in the base-diagram of $(**)$, $v_{U',V'}$ is in $C$. Let $U=u^{-1}(U')$, $V=u^{-1}(V')$, and let us consider the % ?? uncertain: hyphen at line end in "diagramme-base" -- kept as a compound hyphen, or "diagrammebase"
diagram similar to $(**)$, defined by the two inclusions of open sets $U'\subset V'$, $U\subset V$. This time, since the formation of $u_!$ commutes with restrictions to open sets, by transitivity of $\varinjlim$ , the three vertical arrows $c^X_{U'}$, $c^Y_{U'}$, $c^X_{V'}$ are isomorphisms, hence so is $c^f_{U',V'}$. In the same % [sic] vertivales -> verticales, misprint
way, the arrow $Y_V\longrightarrow Y'_{V'}$ (in place of $Y(i)\longrightarrow Y'(i')$) is an isomorphism. Since $v_{U,V}\in C$, the same is true of $v_{U',V'}$, which proves that $f'\in C_{I'^\circ}$. We proceed in the same way for $TC_I$. Finally, the assertion that $u_!$ maps $\mathcal{M}(I^\circ)_{\mathrm{c}}$ into $\mathcal{M}(I'^\circ)_{\mathrm{c}}$ follows from the definitions and from $u_!(C_{I^\circ})\subset C_{I'^\circ}$.

\textbf{Corollary 3.3} (of prop. 3.2 and th. 3.1). \textit{Let $u:I\longrightarrow I'$ be as in \textup{3.2}. Suppose moreover $I$ artinian, and}
\printedpage{28}
% head: A. Grothendieck : Dérivateurs, Ch. XVII, Édition des Universités de Montpellier II et Paris VII
\msp{55}
\textit{that we are for $\mathcal{M}$, $W$, $C$, $I$ under one of the conditions \textup{a), b)} of theorem \textup{3.1}. Then $u_!$ induces an arrow}
\[
u_{!\,\mathrm{cof}}\ :\ \mathcal{M}(I^\circ)_{\mathrm{c}} \longrightarrow \mathcal{M}(I'^\circ)_{\mathrm{c}}\ \subset\ \mathcal{M}(I'^\circ)
\]
\textit{which transforms arrows of $W(I^\circ)$ into arrows of $W(I'^\circ)$ (i.e. compatible with the localizers). It thus induces a functor on the localized categories, whence a}
\[
u_!^W\ :\ \mathrm{Ho}_W(I^\circ) \longrightarrow \mathrm{Ho}_W(I'^\circ) \qquad \mbox{\textsuperscript{(}\footnotemark[17]\textsuperscript{)}}.
\tag*{(3.3.1)}
\]
\footnotetext[17]{\textbf{N.B.} This arrow is defined, even without assuming $I'$ ordered, since $u_{!\,\mathrm{cof}}$ still transforms arrows of $W(I^\circ)$ into arrows of $W(I'^\circ)$. Cf. comment p. 86.}

\textsc{Proof}. By theorem 3.1, we know that $\mathcal{M}(I^\circ)$, $W(I^\circ)$, $C(I^\circ)$ is a structure of a category of models with cofibrations. We also know that $TC_{I^\circ}=C_{I^\circ}\cap W(I^\circ)$, by th. 2.1. $3^\circ$ (p. 27). Thus by 3.2, $u_!$ maps $C_{I^\circ}\cap W(I^\circ)$ into $TC_{I'^\circ}=C_{I'^\circ}\cap W(I'^\circ)$, and a fortiori, into arrows in $W(I'^\circ)$. Moreover, by the theory of categories with cofibrations, we find that every arrow $f$ between cofibrant objects is a composite $f=pi$, where $i\in C_{I^\circ}$, and where $p$ is in $W(I^\circ)$, and moreover a \textit{left inverse} of an arrow in $C_{I^\circ}\cap W(I^\circ)$, i.e. of an arrow of $(TC)_{I^\circ}$. When
\msp{56}
$f\in W(I^\circ)$, then $i\in W(I^\circ)$, hence $i\in W(I^\circ)\cap C_{I^\circ}=(TC)_{I^\circ}$. Altogether, every arrow $f\in W(I^\circ)_{\mathrm{c}}$ is a composite $pi$ in $\mathcal{M}(I^\circ)_{\mathrm{c}}$, where $i\in(TC)_{I^\circ}$, and $p$ is a left inverse of a $j\in(TC)_{I^\circ}$. Since $u_!$ transforms $(TC)_{I^\circ}$ into $(TC)_{I'^\circ}$, it transforms $f$ into $f'=p'i'$, with $i'\in(TC)_{I'^\circ}\subset W_{I'^\circ}$, and $p'$ a left inverse of a $j'=u_!(j)$ which is also in $W_{I'^\circ}$. But then $i'\in W(I'^\circ)$, hence also $f'=p'i'\in W_{I'^\circ}$, which proves that $u_!(W(I^\circ)_{\mathrm{c}})\subset W(I'^\circ)$. The last assertion comes from the following lemma, applied to $(\mathcal{M}(I^\circ),W(I^\circ),C_{I^\circ})$ in place of $(\mathcal{M},W,C)$:

\textbf{Lemma 3.4.} \textit{Let $(\mathcal{M},W,C)$ be a category with cofibrations. Then the inclusion $\mathcal{M}_{\mathrm{c}}\subset\mathcal{M}$ induces an equivalence}
\[
\mathcal{M}_{\mathrm{c}}W_{\mathrm{c}}^{-1} \xrightarrow{\ \sim\ } \mathcal{M}W^{-1}=\mathrm{Ho}_W\ .
\]

We will give the proof (by AQT \texttt{[âne qui trotte: at a donkey's trot, i.e. by routine verification]}) later.

This completes the proof of 3.3.

\textbf{3.5.} The question arises whether the functor $u_!^W$ of (3.3.1) is indeed a left adjoint
\msp{57}
to the functor in the opposite direction $u^*_W$, induced by $u^*$. Let us note that this is the case if $u^*$ transforms $\mathcal{M}(I'^\circ)_{\mathrm{c}}$ into $\mathcal{M}(I^\circ)_{\mathrm{c}}$, and when we assume that $I'$ is also artinian and also satisfies one of the hypotheses a) b) of theorem 3.1. For then $(\mathcal{M}(I^\circ),W(I^\circ),C_{I^\circ})$ is a category with cofibrations by theorem 3.1, we can apply the preceding lemma to it, so $u^*_W$ and $u_!^W$ are realized from a pair of functors
\[
\begin{tikzcd}
\mathcal{M}(I^\circ)_{\mathrm{c}} \arrow[r, shift left, "u_{!}^{\mathrm{c}}"] & \mathcal{M}(I'^\circ)_{\mathrm{c}} \arrow[l, shift left, "u_{\mathrm{c}}^*"]
\end{tikzcd}\ ,
\]